\subsection{The intrinsic image in the joint wild obstruction}

\begin{statement}{E.83}{Proved}{The diagonal obstruction image and empty intrinsic preimages}
Retain the actual maps, source marking and
conjugate paths of E.80. Put
\[
 \begin{gathered}
 N_2=N/\gamma_3N,\qquad N_3=N/\gamma_4N,\qquad
 A_3=\gamma_3N/\gamma_4N=\operatorname{Lie}_3(N^{\rm ab}),\\
 A=\ker(\mathcal J_3\to\mathcal J_2)
       =\mathbb Z_3^4,
       \quad (r,s,r',s')\longleftrightarrow
                     (d^re^s,d^{r'}e^{s'}).
 \end{gathered}
 \tag{E83.1}
\]
The central coefficient map induced by
$N_3\to\mathcal J_3$ factors as the actual
equivariant surjection and inclusion
\[
 \begin{gathered}
 A_3\twoheadrightarrow L\lhook\joinrel\longrightarrow A,\\
 L=\{(r,s,r',s'):r'-r\in3\mathbb Z_3,\quad
                            s+s'-2r\in3\mathbb Z_3\}.
 \end{gathered}
 \tag{E83.2}
\]
In particular these are the image and inclusion
from the integral group computation in E.80.

Use the actual local cohomology coordinates of
E.74 and E.77. They give the natural maps
\[
 \begin{gathered}
 \mathcal B_3=H^2(\mathbb Q_3,A_3)
       \simeq C_3^{63}\oplus C_9^{23}\oplus C_{27}^7,\\
 \mathcal B_3\twoheadrightarrow H^2(\mathbb Q_3,L)
       \longrightarrow H^2(\mathbb Q_3,A).
 \end{gathered}
 \tag{E83.3}
\]
Both groups after $\mathcal B_3$ have the following
fully specified coordinates. A basis for $L$ is
\[
 \begin{gathered}
 (r,s,a,b)\longmapsto
           (r,s,r+3a,2r-s+3b),\\
 K_L=\begin{pmatrix}
 1&0&0&0\\0&1&0&0\\1&0&3&0\\2&-1&0&3
 \end{pmatrix},\qquad \det K_L=9.
 \end{gathered}
 \tag{E83.4}
\]
In this basis and the displayed basis of $A$, the
involutions and Frobenius operators are
\[
 \begin{gathered}
 J_L=\begin{pmatrix}
 -1&0&-3&0\\-2&1&0&-3\\0&0&1&0\\0&0&2&-1
 \end{pmatrix},\qquad
 J_A=-\begin{pmatrix}
 0&0&1&0\\0&0&0&1\\1&0&0&0\\0&1&0&0
 \end{pmatrix},\\
 T_L=T_A=343I_4,\qquad J_AK_L=K_LJ_L.
 \end{gathered}
 \tag{E83.5}
\]
Consequently
\[
 \begin{aligned}
 H^2(\mathbb Q_3,L)&\xrightarrow{\ \sim\ }\mathbb F_3^2,
              &(r,s,a,b)&\longmapsto(r,b-a),\\
 H^2(\mathbb Q_3,A)&\xrightarrow{\ \sim\ }\mathbb F_3^2,
              &(r,s,r',s')&\longmapsto(r+r',s+s'),\\
 H^2(K_L):\mathbb F_3^2&\longrightarrow\mathbb F_3^2,
              &(r,t)&\longmapsto(2r,2r).
 \end{aligned}
 \tag{E83.6}
\]
All coordinates on the right are reduced modulo
three. The image of the full intrinsic map
$\Phi:\mathcal B_3\to H^2(\mathbb Q_3,A)$ is
therefore exactly the diagonal $C_3$. Its exact
remaining quotient is the character
\[
       \mu:\mathbb F_3^2\twoheadrightarrow\mathbb F_3,
                    \qquad \mu(\xi,\eta)=\eta-\xi.
 \tag{E83.7}
\]

Write $\delta_N,\delta_{\mathcal J}$ for the actual
central obstruction maps for $N_3\to N_2$ and
$\mathcal J_3\to\mathcal J_2$. The exact naturality
square is
\[
 \begin{array}{ccc}
 H^1(\mathbb Q_3,N_2)&\longrightarrow&H^1(\mathbb Q_3,\mathcal J_2)\\
 \delta_N\big\downarrow&&\big\downarrow\delta_{\mathcal J}\\
 \mathcal B_3&\xrightarrow{\ \Phi\ }&\mathbb F_3^2.
 \end{array}
 \tag{E83.8}
\]
For the normalized representative of E.80,
\[
       \delta_{\mathcal J}=(\Delta R,\Delta S)\bmod3.
 \tag{E83.9}
\]
Thus the actual intrinsic class-two image is
contained in the explicitly computed locus
$\Delta R\equiv\Delta S\pmod3$.

In particular let $[\sigma_{N,M}]$ be E.80's
descended second-stage section with
$a=y$, $a'=y^{-1}$,
$u=x^{5/2}c^N$, $u'=x^{-1/2}y^2c^M$.
For every $N,M\in\mathbb Z_3$ with
$N+M\not\equiv2\pmod3$ its full intrinsic
class-two preimage is empty:
\[
 \bigl(H^1(\mathbb Q_3,N_2)
          \longrightarrow H^1(\mathbb Q_3,\mathcal J_2)\bigr)^{-1}
                    ([\sigma_{N,M}])=\varnothing.
 \tag{E83.10}
\]
This evaluates a specific obstruction within the
kernel retained in (E80.4), before an intrinsic
class-three lifting question is formed. The
intrinsic preimages of the congruence-compatible
members are retained by that same exact map;
no sufficiency assertion for those preimages is
included in (E83.10).

The whole integral coefficient map is also
evaluated in the retained cubic basis. In the
column order $(f_1,f_2,A_0,A_1,A_2,B_0,B_1,B_2,C_1,C_2)$,
the four geometric unit characters
$(X_h,Y_h,X_{\bar h},Y_{\bar h})$ form the matrix
\[
 Q=\begin{pmatrix}
 2&-1&0&0&0&0&0&0&0&-3\\
 1&0&0&0&0&1&1&1&0&-3\\
 1&-2&0&0&0&0&0&0&-3&0\\
 0&-1&0&0&0&1&1&1&-3&0
 \end{pmatrix}.
 \tag{E83.11}
\]
For a degree-three Hall word $[[q_a,q_b],q_c]$,
the actual $4$ by $330$ matrix $F_3:A_3\to A$
has the following two coordinate pairs, using
$\epsilon=0,2$ for the two pairs of rows of $Q$:
\[
 -\bigl(Q_{\epsilon,a}Q_{\epsilon+1,b}
         -Q_{\epsilon+1,a}Q_{\epsilon,b}\bigr)
           (Q_{\epsilon,c},Q_{\epsilon+1,c}).
 \tag{E83.12}
\]
Here row and column indices of $Q$ start at zero; the Hall
ordering is exactly E.41's retained ordering.
The complete matrix is saved, together with its
actual arithmetic intertwiners and image basis.

Let $b_1,\ldots,b_{93}$ be E.74's positive Smith
coordinates, in their retained order. The first
$63$ have order three, the next $23$ have order
nine and the final seven have order twenty-seven.
Then the full intrinsic obstruction map and its
kernel have the explicit coordinates
\[
 \begin{gathered}
 \Phi(b)=(\lambda(b),\lambda(b)),\qquad
 \lambda(b)=b_{89}+b_{91}+2b_{92}\pmod3,\\
 \ker\Phi\simeq C_3^{63}\oplus C_9^{24}\oplus C_{27}^6,\\
 k=(b_{89}+b_{91}+2b_{92})/3\in C_9,\qquad
 b_{89}=3k-b_{91}-2b_{92}\in C_{27}.
 \end{gathered}
 \tag{E83.13}
\]
For these inverse kernel coordinates retain all
$b_j$ except $b_{89}$ and replace it by $k$.
Division by three in the displayed subgroup of
$C_{27}$ means its unique inverse identification
with $C_9$.

Let $\mathscr I$ be E.82's sum of the nine actual
cusp-family obstruction images. The entire
subgroup $\mathscr I$ lies in $\ker\Phi$.
In E.82's retained residual coordinates
$(\eta_1,\ldots,\eta_{12})$ on
$\mathcal B_3/\mathscr I=C_3^{11}\oplus C_9$,
the exact induced map and its kernel are
\[
 \begin{gathered}
 \bar\lambda(\eta)=\eta_9+\eta_{10}\pmod3,\\
 (\ker\Phi)/\mathscr I\simeq C_3^{10}\oplus C_9,
       \qquad \eta_9=-\eta_{10}.
 \end{gathered}
 \tag{E83.14}
\]
Its inverse retains the other eleven residual
coordinates and inserts the displayed ninth one.
Thus the cohomological comparison also computes
one exact quotient of the remaining E.82
obstruction coordinates.

\emph{Manuscript: \S51.16, folios733v--736v,
scans172--178, formulas(371)--(379), Part E(c),(d)
and follow-up(4). The result computes the exact
cohomological image of the integral central map
and its effect on actual descended quotient
sections.}
\end{statement}

\begin{proof}
The maps in (E83.2) are E.80's actual central maps.
Their image is $L=\gamma_3\mathcal J_3$, since
$N_3\to\mathcal J_3$ is onto. The lattice is
free of rank four; it has the basis (E83.4),
obtained by writing $r'=r+3a$ and
$s'=2r-s+3b$. The first map is onto this lattice
as an integral module. Its kernel is a saturated
free submodule of $A_3$, because the quotient
$L$ is torsion-free. The inclusion $L\to A$
has the separately retained index nine.

On a pure third-degree target coordinate, both
actual E.80 operators of $\tau$ and $\tau'$ are
multiplication by $7^3$. The involution is minus
the factor swap, by (E80.19). This gives $T_A,J_A$.
Substitution of the basis (E83.4) gives the integer
matrix $J_L=K_L^{-1}J_AK_L$ displayed in (E83.5).
In particular $J_L^2=I$; no rational identification
of $L$ with all of $A$ is used.

The local-duality proof in (E74.15) identifies
$H^2(\mathbb Q_3,M)$ naturally with the compact
coinvariants of $M(-1)$ for each of these finite
free modules. Coefficient homomorphisms preserve
the finite-level evaluation pairing, and the
inverse-limit identifications are compatible with
them. Equivalently E.75's actual relative source
resolution gives the same coefficient maps. Since
the local image is the actual $\Gamma$, this
natural presentation is
\[
 H^2(\mathbb Q_3,M)
     =M/\bigl((T_M-7I)M+(J_M+I)M\bigr),
             \qquad M=A_3,L,A.
 \tag{E83.15}
\]
The same proof applies to $L$ and $A$: their
actions factor through that actual quotient, and
all lattices and finite coefficient systems are
finite free before reduction. The images in the
denominator are compact. This also proves the
surjectivity of the first arrow in
(E83.3) directly: lift a representative in $L$
through the coefficient surjection $A_3\to L$
and pass to (E83.15).

For $M=A$ or $L$, $T_M-7I=336I=3\cdot112I$,
and $112$ is a three-adic unit. Reduce the
involution relations modulo three. For $A$ they
identify the two factors; the complete quotient
coordinates are $(r+r',s+s')$. For $L$, the
image of $J_L+I$ modulo three is spanned by
the $s$ basis vector and the sum of the $a,b$
basis vectors. Its quotient coordinates are
$(r,b-a)$. These observations prove both
isomorphisms in (E83.6), including their
surjectivity and exact kernels. The actual
inclusion sends
\[
  (r,s,a,b)\longmapsto(r,s,r+3a,2r-s+3b)
   \longmapsto(2r,2r)\pmod3.
 \tag{E83.16}
\]
This proves the displayed natural map, its
diagonal image, and the residual quotient (E83.7).
In particular the intermediate exact sequence is
$0\to\mathbb F_3\to H^2(\mathbb Q_3,L)
  \to H^2(\mathbb Q_3,A)\to\mathbb F_3\to0$,
with first map $t\mapsto(0,t)$ and last map $\mu$.

There is also an explicit nonzero image witness
in the actual intrinsic lattice. Choose geometric
lifts $n_1,n_2\in N$ of $g_1,g_2$ in (E80.23),
using the proved geometric surjection. The word
$[n_2,[n_1,n_2]]$ maps to the central vector
$(1,1,1,1)$; this is the third column of (E80.26).
Its cohomological image is $(2,2)$, a generator
of the diagonal. This is consistent with, and
also witnesses, the surjectivity argument above.

To prove the obstruction naturality with its
normalization, use the literal equivariant
diagram of central extensions
\[
 \begin{array}{ccccccccc}
 1&\longrightarrow&A_3&\longrightarrow&N_3&\longrightarrow&N_2&\longrightarrow&1\\
 &&\big\downarrow&&\big\downarrow&&\big\downarrow\\
 1&\longrightarrow&A&\longrightarrow&\mathcal J_3
       &\longrightarrow&\mathcal J_2&\longrightarrow&1.
 \end{array}
 \tag{E83.17}
\]
The left arrow is the composite (E83.2), rather
than a surjection onto the full bottom kernel.
Choose continuous Hall lifts of a lower cocycle
to $N_3$ and apply the middle arrow. Its
multiplication defect is the image of the
$A_3$-valued defect. Changing to the target's
chosen Hall lifts adds an $A$-valued coboundary.
Centrality makes the coefficient action independent
of the lower geometric conjugator. This proves
(E83.8) on the full section sets. It also proves
that its assertion is independent of the retained
representatives and gauges.

In E.80's normalized source notation let
$v_-=v_{-5/2}(u,a)$ and
$v_+=v_{1/2}(u',a')$. The second-stage relation
makes $v_-v_+^{-1}$ central in $C_3$, with
coordinates $\Delta=(\Delta R,\Delta S)$.
The joint lifted word is $(v_-,Iv_+)$.
Since $I$ is multiplication by minus one in
third degree, its anti-invariant source defect
and the E.77 invariant are exactly
\[
 \omega=(\Delta,\Delta),\qquad
       [\omega/2]=[(\Delta/2,\Delta/2)]\in H^2(\mathbb Q_3,A).
 \tag{E83.18}
\]
Applying the target coordinate sum gives precisely
$\Delta\bmod3$, which proves (E83.9) without a
factor-of-two ambiguity. The map $\mu$ vanishes
on the image of $\Phi$, so (E83.8) proves the
claimed necessary locus for every actual
intrinsic class-two section.

For the family in (E80.14), the computed target
defect is $(0,7(5-N-M))$. Its residual value is
$7(5-N-M)\bmod3$, which is nonzero exactly
for the parameters in (E83.10). An intrinsic
class-two preimage would contradict the exact
commutative square and the diagonal image.
This proves the empty-preimage assertion.

\textbf{The full coefficient and residual maps.}
The positive peripheral coefficient of a unit
character is its order at that cusp. E.54 gives
$\operatorname{div}(h)=3O-3C_2$ and
$\operatorname{div}(1-h)=\sum_jB_j-3C_2$.
Conjugating swaps $C_1,C_2$. These are exactly
the last eight columns of (E83.11), in E.37's
positive peripheral convention. Its actual
integral gluing has columns $j_*=[3I;U]$, with
$U_d=(v_d,-u_d)$ in that section's table.
The composite of the unit homomorphism from
the generalized Jacobian with $j:E\to J_{E,D}$
is trivial: a group morphism from the proper
connected elliptic curve to $\mathbb G_m$ is
constant, and its value at the identity is one.
Thus a unit character with toric row $m_T$
has its exact elliptic row determined by
\[
                3m_f+m_TU=0,
                    \qquad m_f=-m_TU/3.
 \tag{E83.19}
\]
This proves all four first blocks of (E83.11)
integrally, with their signs. The rows for
$\ell$, $X_h$, $(X_h+X_{\bar h})/3$ have an
integer determinant-one minor in columns
$(f_2,B_0,C_1)$. They therefore also recover
the surjection onto E.80's joint abelian lattice.

The precise E.78 marking transport also preserves
the retained coefficient matrices. It replaces
$\tau_0$ by $h_0\tau_0h_0^{-1}$. Over $K$ the
E.37 matrix on $V$ is
$\left(\begin{smallmatrix}A_E&0\\
(UA_E-\chi U)/3&\chi I\end{smallmatrix}\right)$.
Over $\mathbb Q_3$ this is conjugate, by the fixed
matrix $[3I,0;U,I]$, to
$\operatorname{diag}(A_E,\chi I)$. The CM matrices
$A_E$ commute, and therefore these actual integral
matrices commute as well. Conjugation by the
matrix of $h_0$ leaves $T_1$ unchanged; the
arithmetic involution is the retained one.
Thus $T_3,J_3$ in the full coefficient calculation
are exactly E.74's matrices after the actual
source-word transport, with no change of the
coefficient coordinate system.

A homomorphism of free pro-three groups induces
the free Lie homomorphism on its lower central
graded groups. In the target convention
$c=[x,y]$, $d=[x,c]$, $e=[y,c]$, a second-degree
determinant is the coefficient of $c$, and
$[c,Ax+By]=-Ad-Be$. This proves (E83.12).
The resulting image has the basis of (E83.4)
up to an integer unimodular change, as verified
by its determinant-nine Hermite basis. The
arithmetic equalities are the exact identities
$F_3T_3=343F_3$ and $F_3J_3=J_AF_3$.

Let $L_*$ and $L_*^{-1}$ be E.74's retained
left Smith matrix and inverse modulo $729$,
and let $I_+$ be its ordered set of positive
rows. The complete coordinate map on the
source obstruction group is
\[
 \begin{pmatrix}1&0&1&0\\0&1&0&1\end{pmatrix}
      F_3(L_*^{-1})_{*,I_+}\pmod3.
 \tag{E83.20}
\]
The columns belonging to the omitted unit
rows vanish modulo three. Both rows of
(E83.20) equal the character in (E83.13).
The three nonzero entries occur at the
specified order-twenty-seven coordinates,
with coefficients $1,1,2$. These statements
are exact modulo three, so the retained
inverse precision gives the actual map.
The formula for $k$ and its inverse then
proves the kernel type and both coordinate
maps in (E83.13).

There is a geometric reason that every E.82
cusp image is killed. Modulo three, (E83.11)
sends each $B_j$ to $(y,y)$ and every other
cusp cocharacter to zero. A second-degree
pair has coordinates $(n,n')$ with
$n+n'=0$ modulo three, by (E80.15).
Its bracket with $(y,y)$ has central
coordinates $(0,-n,0,-n')$, whose image
under the target sum in (E83.6) is zero.
This proves $\Phi\delta_i=0$ for all nine
actual maps, hence $\mathscr I\subset\ker\Phi$.

For the numerical residual coordinate map,
let $P_*$ be E.82's exact left Smith matrix
of its $93$ by $93$ relation lattice and
let $I'_+$ be its twelve positive positions.
Using the full row $\lambda$ in (E83.13),
the quotient character is
\[
             \lambda(P_*^{-1})_{*,I'_+}\pmod3
                 =(0,0,0,0,0,0,0,0,1,1,0,0).
 \tag{E83.21}
\]
All rational denominators in this inverse
matrix are prime to three. The unit positions
vanish, and all relation-lattice columns are
killed, so the displayed row is the literal
induced homomorphism. This proves (E83.14),
including its surjectivity, kernel and inverse.

The self-contained program
\texttt{a\_joint\_intrinsic\_image.py} reconstructs
the actual Hall image matrix from the two
integral abelian lattices. It saves its integer
Smith maps and inverses, those for both local
cohomology relation matrices, the integral
involution on $L$, and the explicit natural
cohomology map. It verifies all $81$ ambient
residue vectors, all $729$ inputs to the six
actual Hall image columns, all nine index-two
normalizations, and all $81$ family residues
with parameters modulo nine. Every check passes;
the program and its actual output and JSON
are retained.
The second program
\texttt{a\_joint\_intrinsic\_coordinates.py}
derives the four unit characters, checks all
$660$ target tensor expansions independently,
reconstructs the entire $4$ by $330$ map,
checks both exact arithmetic intertwiners,
and saves both actual cohomology matrices.
It verifies the vanishing of all $405$
E.82 cusp columns and computes the complete
residual row and its coordinate inverse.
The $6561$ triples for the changed kernel
coordinates in (E83.13) and all three values
of the residual coordinate in (E83.14) are
also checked in both directions.
Its retained input hashes specify the exact
E.37/E.41, E.74 and E.82 coordinate systems
used in (E83.13)--(E83.14).
\end{proof}
